\textbf{PROPOSAL TUGAS AKHIR}

\textbf{JUDUL TUGAS AKHIR MAKSIMAL 20 KATA}

\textbf{Oleh:}

\textbf{NAMA MAHASISWA KAPITAL SEMUA}

\textbf{NIM Mahasiswa dalam angka}

\includegraphics[width=1.5367in,height=2.22917in]{docxmedia/media/image1.jpg}

\textbf{PROGRAM STUDI STRATA 1 SISTEM INFORMASI}

\textbf{FAKULTAS REKAYASA INDUSTRI}

\textbf{UNIVERSITAS TELKOM}

\textbf{202\hl{x}}

\hypertarget{daftar-isi}{%
\section{DAFTAR ISI}\label{daftar-isi}}

ABSTRAK \protect\hyperlink{_Toc211426824}{ii}

\emph{ABSTRACT} \protect\hyperlink{_Toc211426825}{iii}

LEMBAR PENGESAHAN \protect\hyperlink{_Toc211426826}{iv}

LEMBAR PERNYATAAN ORISINALITAS \protect\hyperlink{_Toc211426827}{v}

KATA PENGANTAR \protect\hyperlink{_Toc211426828}{vi}

DAFTAR ISI \protect\hyperlink{daftar-isi}{vii}

DAFTAR GAMBAR \protect\hyperlink{daftar-gambar}{ix}

DAFTAR TABEL \protect\hyperlink{daftar-tabel}{x}

DAFTAR ISTILAH \protect\hyperlink{daftar-istilah}{xi}

DAFTAR SINGKATAN \protect\hyperlink{daftar-singkatan}{xii}

DAFTAR SIMBOL \protect\hyperlink{daftar-simbol}{xiii}

DAFTAR LAMPIRAN \protect\hyperlink{daftar-lampiran}{xiv}

Bab I PENDAHULUAN \protect\hyperlink{pendahuluan}{1}

I.1 Latar Belakang \protect\hyperlink{latar-belakang}{1}

I.2 Rumusan Masalah \protect\hyperlink{rumusan-masalah}{4}

I.3 Tujuan Penelitian \protect\hyperlink{tujuan-penelitian}{5}

I.4 Batasan Penelitian \protect\hyperlink{batasan-penelitian}{5}

I.5 Manfaat Penelitian \protect\hyperlink{manfaat-penelitian}{6}

I.6 Sistematika Penulisan \protect\hyperlink{sistematika-penulisan}{6}

Bab II TINJAUAN PUSTAKA \protect\hyperlink{tinjauan-pustaka}{8}

II.1 Materi-materi pendukung riset
\protect\hyperlink{materi-materi-pendukung-riset-judul-sub-bab-silakan-disesuaikan-tidak-harus-seperti-ini}{8}

II.2 Penelitian Terdahulu (Judul Sub bab tidak harus seperti ini,
silakan sesuaikan)
\protect\hyperlink{penelitian-terdahulu-judul-sub-bab-tidak-harus-seperti-ini-silakan-sesuaikan}{13}

Bab III METODE PENYELESAIAN MASALAH
\protect\hyperlink{metode-penyelesaian-masalah}{16}

III.1 Metode Penelitian \protect\hyperlink{metode-penelitian}{16}

III.2 Sistematika Penyelesaian Masalah
\protect\hyperlink{sistematika-penyelesaian-masalah}{16}

III.2.1 Tahap Pertama \protect\hyperlink{tahap-pertama}{17}

III.2.2 Tahap Kedua \protect\hyperlink{tahap-kedua}{17}

III.3 Metode Pengumpulan data dan Pengolahan Data (Opsional)
\protect\hyperlink{metode-pengumpulan-data-dan-pengolahan-data-opsional}{17}

III.4 Metode Evaluasi \protect\hyperlink{metode-evaluasi}{17}

Bab IV Hasil PENELITIAN \protect\hyperlink{_Toc211426853}{19}

IV.1 Identifikasi Kebutuhan Data dan Analisa
\protect\hyperlink{_Toc211426854}{19}

IV.1.1 Identifikasi Kebutuhan \protect\hyperlink{_Toc211426855}{19}

IV.1.2 Sub Sub Bab Selanjutnya \protect\hyperlink{_Toc211426856}{19}

IV.2 Arsitektur atau Model Konseptual (Opsional)
\protect\hyperlink{_Toc211426857}{19}

IV.3 Perancangan Sistem \protect\hyperlink{_Toc211426858}{21}

IV.4 Hasil Implementasi Sistem \protect\hyperlink{_Toc211426859}{21}

Bab V VALIDASI HASIL DAN DISKUSI \protect\hyperlink{_Toc211426860}{22}

V.1 Evaluasi Penelitian \protect\hyperlink{_Toc211426861}{22}

V.2 Diskusi \protect\hyperlink{_Toc211426862}{24}

V.3 Dampak Hasil Penelitian \protect\hyperlink{_Toc211426863}{24}

Bab VI KESIMPULAN DAN SARAN \protect\hyperlink{_Toc211426864}{25}

VI.1 Kesimpulan \protect\hyperlink{_Toc211426865}{25}

VI.2 Saran \protect\hyperlink{_Toc211426866}{25}

DAFTAR PUSTAKA \protect\hyperlink{daftar-pustaka}{26}

LAMPIRAN \protect\hyperlink{lampiran}{27}

\hypertarget{daftar-gambar}{%
\section{DAFTAR GAMBAR}\label{daftar-gambar}}

\protect\hyperlink{_Toc211422568}{Gambar I‑1 Contoh Tulisan Paragraf
dengan Ruang Kosong \protect\hyperlink{_Toc211422568}{4}}

\protect\hyperlink{_Toc211422569}{Gambar II‑1 Tampilan halaman
\emph{dashboard} monitoring \protect\hyperlink{_Toc211422569}{9}}

\protect\hyperlink{_Toc211422570}{Gambar II‑2 Perbandingan tingkat
penjualan dan jumlah produk \protect\hyperlink{_Toc211422570}{10}}

\protect\hyperlink{_Toc211422571}{Gambar III‑1 Alur sistematika
penelitian \protect\hyperlink{_Toc211422571}{17}}

\protect\hyperlink{_Toc211422572}{Gambar IV‑1 Contoh bentuk arsitektur
sistem yang diusulkan \protect\hyperlink{_Toc211422572}{20}}

\protect\hyperlink{_Toc211422573}{Gambar IV‑2 Contoh bentuk arsitektur
pada bidang data \protect\hyperlink{_Toc211422573}{21}}

\protect\hyperlink{_Toc211422574}{Gambar V‑1 Contoh hasil penelitian
dalam grafik \protect\hyperlink{_Toc211422574}{23}}

\protect\hyperlink{_Toc211422575}{Gambar V‑2 Contoh hasil penelitian
dalam grafik \protect\hyperlink{_Toc211422575}{23}}

\hypertarget{daftar-tabel}{%
\section{DAFTAR TABEL}\label{daftar-tabel}}

\protect\hyperlink{_Toc211422712}{Tabel II‑1 Perbancingan contoh tugas
\protect\hyperlink{_Toc211422712}{11}}

\protect\hyperlink{_Toc211422713}{Tabel II‑2 Posisi Penelitian
(\emph{State of The Art)} atau judul dapat disesuaikan dengan konteksnya
\protect\hyperlink{_Toc211422713}{15}}

\protect\hyperlink{_Toc211422714}{Tabel V‑1 Contoh tabel hasil
perbandingan metode untuk evaluasi
\protect\hyperlink{_Toc211422714}{22}}

\hypertarget{daftar-istilah}{%
\section{DAFTAR ISTILAH}\label{daftar-istilah}}

\begin{longtable}[]{@{}
  >{\raggedright\arraybackslash}p{(\columnwidth - 4\tabcolsep) * \real{0.2179}}
  >{\raggedright\arraybackslash}p{(\columnwidth - 4\tabcolsep) * \real{0.5281}}
  >{\raggedright\arraybackslash}p{(\columnwidth - 4\tabcolsep) * \real{0.2541}}@{}}
\toprule\noalign{}
\begin{minipage}[b]{\linewidth}\raggedright
Istilah
\end{minipage} & \begin{minipage}[b]{\linewidth}\raggedright
Deskripsi
\end{minipage} & \begin{minipage}[b]{\linewidth}\raggedright
Halaman Pertama Kali Digunakan
\end{minipage} \\
\midrule\noalign{}
\endhead
\bottomrule\noalign{}
\endlastfoot
AI & Singkatan dari \emph{Artificial Intelligence} yang merujuk pada
kemampuan sistem komputer untuk menjalankan tugas-tugas yang biasanya
membutuhkan kecerdasan manusia. & 2 \\
Scrum & Scrum adalah \ldots{} & 4 \\
Istilah & Jelaskan istilah-istilah yang pembaca pada umumnya tidak tahu.
Untuk membantu jika pembaca tidak paham. Silakan pilih kosakata istilah
yang tidak umum. Urutkan istilah secara abjad, serta di halaman mana
pertama kali disebutkan/digunakan & 5 \\
& & \\
\end{longtable}

\hypertarget{daftar-singkatan}{%
\section{DAFTAR SINGKATAN}\label{daftar-singkatan}}

\begin{longtable}[]{@{}
  >{\raggedright\arraybackslash}p{(\columnwidth - 4\tabcolsep) * \real{0.2179}}
  >{\raggedright\arraybackslash}p{(\columnwidth - 4\tabcolsep) * \real{0.5281}}
  >{\raggedright\arraybackslash}p{(\columnwidth - 4\tabcolsep) * \real{0.2541}}@{}}
\toprule\noalign{}
\begin{minipage}[b]{\linewidth}\raggedright
Singkatan
\end{minipage} & \begin{minipage}[b]{\linewidth}\raggedright
Deskripsi
\end{minipage} & \begin{minipage}[b]{\linewidth}\raggedright
Halaman Pertama Kali Digunakan
\end{minipage} \\
\midrule\noalign{}
\endhead
\bottomrule\noalign{}
\endlastfoot
IoT & \emph{Internet of Things} & 2 \\
KDD & \emph{Knowledge Discovery in Databases} & 4 \\
& & \\
& & \\
\end{longtable}

\hypertarget{daftar-simbol}{%
\section{DAFTAR SIMBOL}\label{daftar-simbol}}

\begin{longtable}[]{@{}
  >{\raggedright\arraybackslash}p{(\columnwidth - 4\tabcolsep) * \real{0.2179}}
  >{\raggedright\arraybackslash}p{(\columnwidth - 4\tabcolsep) * \real{0.5281}}
  >{\raggedright\arraybackslash}p{(\columnwidth - 4\tabcolsep) * \real{0.2541}}@{}}
\toprule\noalign{}
\begin{minipage}[b]{\linewidth}\raggedright
Simbol
\end{minipage} & \begin{minipage}[b]{\linewidth}\raggedright
Deskripsi
\end{minipage} & \begin{minipage}[b]{\linewidth}\raggedright
Halaman Pertama Kali Digunakan
\end{minipage} \\
\midrule\noalign{}
\endhead
\bottomrule\noalign{}
\endlastfoot
\(\sigma\) & Simpangan baku atau standar deviasi & 5 \\
\(\mu\) & Nilai rata-rata & 10 \\
& & \\
\end{longtable}

\hypertarget{daftar-lampiran}{%
\section{DAFTAR LAMPIRAN}\label{daftar-lampiran}}

\protect\hyperlink{_Toc211427018}{LAMPIRAN A JUDUL LAMPIRAN
\protect\hyperlink{_Toc211427018}{27}}

\protect\hyperlink{_Toc211427019}{LAMPIRAN B DOKUMENTASI WAWANCARA
\protect\hyperlink{_Toc211427019}{28}}

\hypertarget{pendahuluan}{%
\section{PENDAHULUAN}\label{pendahuluan}}

Bab Pendahuluan ini menyajikan kerangka dasar penelitian, mulai dari
latar belakang masalah, perumusan masalah dan tujuan, batasan, manfaat
studi, hingga diakhiri dengan sistematika penulisan sebagai panduan
untuk keseluruhan laporan.

\hypertarget{latar-belakang}{%
\subsection{\texorpdfstring{Latar Belakang
}{Latar Belakang }}\label{latar-belakang}}

Latar Belakang menyajikan deskripsi area topik penelitian serta isu atau
problematika yang saat ini membutuhkan solusi pemecahan. Subbab ini
diawali dengan deskripsi topik kajian yang didukung dengan data,
informasi, atau referensi yang dapat memperlihatkan pentingnya
penelitian dilakukan. Data yang disajikan dapat menunjukan permasalahan
(contoh: keterlambatan, kesalahan, dan lain-lain) atau data yang
menunjukan adanya kekurangan (gap) antara kondisi eksisting dengan
target (\emph{needs}) yang akan dicapai.

Setelah menampilkan permasalahan yang didukung dengan data, mahasiswa
harus mampu menganalisis akar permasalahan dari permasalahan atau
identifikasi gap yang telah ditampilkan. Identifikasi akar pemasalahan
atau gap tersebut menjadi dasar perumusan masalah yang dituliskan di
subbab selanjutnya.

Dalam melakukan analisis akar permasalahan atau gap, mahasiswa
diharapkan dapat mengaitkan permasalahan tersebut dengan kondisi
organisasi atau \emph{stakeholders} yang terdampak. Mahasiswa juga harus
mampu menjelaskan secara jelas nilai (\emph{value}) atau dampak positif
yang akan diperoleh organisasi (seperti efisiensi proses, peningkatan
kualitas keputusan, efisiensi biaya, atau efektivitas layanan) apabila
permasalahan tersebut diselesaikan melalui penelitian/sistem yang
dikembangkan

Pada Gambar I-1 adalah gambar terkait dengan gambar contoh. Silakan
sebutkan nomor gambar terlebih dahulu baru dimasukkan gambarnya. Seluruh
gambar harus dirujuk pada teks.

\includegraphics[width=4.9558in,height=2.55451in]{docxmedia/media/image2.png}

Gambar I-1. Harga jual cabai rawit merah di beberapa provinsi

Sumber: Badan Pangan Nasional (2024)

Untuk mendukung pernyataan terkait permasalahan, maka diwajibkan
menggunakan dukungan literatur. Misal jika menyebutkan pertumbuhan
pengguna aplikasi ERP meningkat 20\% per tahun sebutkan selalu
referensinya (Reffad \& Alti, 2023). Demikian pula jika menyebutkan
permasalahan-permasalahan yang diambil atau disarikan dari publikasi
penelitian sebelumnya, selalu sebutkan referensinya. (Abdelaziz dkk.,
2021) . Sitasi yang diberikan di sini hanya contoh saja bagaimana setiap
kalimat yang bukan pernyataan dari penulis, harus didukung dengan
literatur. Contohnya adalah sebagai berikut. Menurut Suakanto dkk.
(2021) manajemen asset telah memiliki siklus yang jelas dari awal dan
akhir. Pernyataan ini tentunya bukan pernyataan dari penulis karena hal
ini sudah banyak dituliskan oleh peneliti yang lain. Penulis harus
menggunakan dukungan literatur terhadap pernyatan-pernyataan tersebut.

Pastikan ada relevansi dengan Subbab I.2 dan I.3. Jika merumuskan
masalah di subbab I.2, maka hal tersebut juga telah dijelaskan terlebih
dahulu secara umum di bagian Latar Belakang ini. Usahakan jangan
tiba-tiba muncul rumusan masalah, tetapi tidak pernah dibahas di bagian
latar belakang. Sehingga cerita dari awal sampai di rumusan masalah
terkesan tidak menyambung.

Usahakan menulis secara mengalir. Sehingga jika ingin berpindah dari
satu ide ke ide pokok yang lain, maka harus membuat \emph{bridge}.
\emph{Bridge} dapat berupa kalimat untuk menghubungkan antara satu
paragraf dengan paragraf yang lain. Sehingga alur cerita dari satu
paragraf ke paragraf yang lain bisa menyambung dan mengalir.

Dalam penyusunan tugas akhir, mahasiswa diperbolehkan menggunakan
teknologi \emph{Artificial Intelligence} (AI) sebagai alat bantu
pendukung proses penelitian. Meskipun demikian, mahasiswa tidak boleh
melakukan salin-tempel langsung (\emph{copy-paste}) dari hasil keluaran
AI ke dalam naskah. Mahasiswa perlu melakukan penyuntingan dan penulisan
ulang (\emph{paraphrasing/rewriting}) menggunakan gaya bahasa, analisis,
dan pemikiran sendiri untuk menjaga originalitas karya ilmiah.

Sebagai bentuk transparansi etika akademik, mahasiswa wajib melampirkan
dokumentasi penggunaan AI pada bagian Lampiran. Dokumentasi ini berupa
\emph{screenshot} (tangkapan layar) yang memperlihatkan secara jelas
instruksi (\emph{prompting}) yang diajukan beserta hasil keluaran
(\emph{output}) dari AI. Hal ini bertujuan untuk menunjukkan secara
faktual bagaimana teknologi AI dimanfaatkan secara proporsional dalam
mendukung penyelesaian tugas akhir ini.

Mahasiswa wajib bertanggung jawab penuh terhadap seluruh isi naskah dan
harus mampu menjelaskan serta mempertahankan setiap bagian yang
dituliskan saat proses \emph{review} atau sidang proposal. Penggunaan AI
yang berlebihan hingga menyebabkan mahasiswa tidak memahami isi
tulisannya sendiri akan berdampak langsung pada penilaian. Apabila pada
saat \emph{review} mahasiswa terbukti tidak mampu menjelaskan materi
akibat terlalu bergantung pada AI, maka penguji/reviewer memiliki hak
untuk tidak meluluskan proposal tersebut atau mewajibkan mahasiswa
melakukan presentasi ulang.

Aturan tata tulis layout mengikuti aturan dari Panduan Tugas Akhir 2024.
Secara umum adalah Naskah dicetak pada kertas HVS berukuran A4 (210 mm x
297 mm) dan berat 80 g/m2 (HVS 80 GSM). Batas cetak adalah 4 cm dari
tepi kiri kertas, 3 cm dari tepi kanan kertas, dan 3 cm dari tepi atas
dan tepi bawah kertas. Awal paragraf sejajar dengan penomoran subbab.
Baris-baris kalimat naskah berjarak satu setengah spasi. Baris notasi
pada gambar, tabel, judul tabel, judul gambar, dan daftar pustaka
berjarak satu spasi. Antar paragraf berjarak sesuai default dari Style
Normal (''add space between paragraphs''). Awal paragraf tidak
diindentasi. Huruf pertama paragraf baru dimulai dari batas tepi kiri
naskah. Tidak memulai paragraf baru pada dasar halaman, kecuali apabila
cukup tempat untuk sedikitnya dua baris. Huruf pertama sesudah
tanda-baca koma (,), titik-koma (;), titik-ganda (:) dan titik (.)
dicetak dengan menyisihkan suatu rongak (ruangan antara dua huruf) di
belakang tanda-baca tersebut. Nomor halaman berkelanjutan.

Baris antar kalimat pada paragraf dalam laporan Tugas Akhir berjarak
satu setengah spasi (1,5 line) dengan jarak tambahan antar paragraf 8pt.
Jarak satu spasi (1 line) digunakan pada catatan kaki, judul keterangan
dengan isi tabel, tabel, gambar, dan daftar pustaka. Huruf pertama pada
paragraf baru dimulai dari tepi batas kiri naskah. Paragraf tidak boleh
dimulai di dasar halaman, kecuali terdapat minimal dua baris. Baris
terakhir dari sebuah paragraf tidak boleh berada pada awal halaman baru
berikutnya, upayakan baris terakhir paragraf terdapat pada dasar halaman
sebelumnya. Huruf pertama sesudah tanda baca koma (,), titik-koma (;),
titik-ganda (:), titik (.), tanda tanya (?), tanda seru (!), dan tanda
baca lainnya dicetak dengan menyisakan satu rongak (ruangan antara dua
huruf) di belakang tanda baca tersebut. Penulisan bahasa dan istilah
yang baku mengikuti aturan dan kelaziman penulisan dalam disiplin
keilmuan yang diikuti. Setiap bab baru dimulai dengan halaman baru. Satu
paragraf minimal terdiri dari dua kalimat. Judul Subbab merupakan bagian
besar dari Bab dan minimal terdiri dari 2 paragraf.

Bab baru diawali di halaman baru. Baris-baris kalimat naskah tugas akhir
berjarak satu setengah spasi. Penulisan paragraf menggunakan Styles
``Normal'', yaitu font Times New Roman ukuran 12pt. Dalam penulisan,
upayakan sebuah halaman diisi dengan optimal. Jangan biarkan ada ruang
kosong cukup besar misal karena memasukkan sebuah gambar atau foto.
Ruang kosong tersebut perlu diisi dengan teks yang ada sesudah gambar
tersebut. Ilustrasinya dapat dilihat pada Gambar I-2.

\includegraphics[width=5.51181in,height=3.72153in]{docxmedia/media/image3.png}

\protect\hypertarget{_Toc211422568}{}{}Gambar I‑1 Contoh Tulisan
Paragraf dengan Ruang Kosong

\hypertarget{rumusan-masalah}{%
\subsection{\texorpdfstring{Rumusan Masalah
}{Rumusan Masalah }}\label{rumusan-masalah}}

Rumusan masalah yang mendasari penelitian ini adalah:

\begin{enumerate}
\def\labelenumi{\alph{enumi}.}
\item
  Rumusan masalah pertama
\item
  Rumusan masalah kedua. Rumusan masalah disusun dalam dua kalimat.
  Kalimat pertama adalah kalimat pernyataan yang merupakan problem yang
  dihadapi seperti apa. Selanjutnya diikuti dengan pertanyaan
  penelitiannya.
\item
  Rumusan masalah ketiga. Sebagai contoh adalah sebagai berikut.
  Penyakit pada tanaman teh susah dideteksi dengan menggunakan mata
  manusia biasa. Apakah penggunaan model deteksi dari kamera smartphone
  menggunakan model CNN dapat mempredeksi dengan akurat pada penyakit
  tanaman teh?
\item
  Rumusan masalah keempat. Sebagai contoh adalah sebagai berikut. Proses
  pengajuan surat menyurat di kelurahan memakan waktu yang lama yang
  biasanya berkisar 3 hari sampai seminggu. Apakah pengembangan aplikasi
  web untuk surat menyurat di kelurahan dapat membantu efisiensi
  pembuatan surat menyurat hanya kurang dalam 1 hari saja?
\end{enumerate}

\hypertarget{tujuan-penelitian}{%
\subsection{Tujuan Penelitian}\label{tujuan-penelitian}}

Berdasarkan rumusan masalah sebelumnya, tujuan penelitian dapat
diuraikan sebagai berikut:

\begin{enumerate}
\def\labelenumi{\alph{enumi}.}
\item
  Tujuan penelitian pertama -- usahakan merupakan tujuan untuk menjawab
  rumusan masalah pertama dan seterusnya
\item
  Tujuan penelitian kedua -- menjawab rumusan masalah kedua
\item
  Tujuan penelitian ketiga
\item
  Tujuan penelitian keempat
\end{enumerate}

\hypertarget{batasan-penelitian}{%
\subsection{Batasan Penelitian}\label{batasan-penelitian}}

Penelitian ini dibatasi pada beberapa hal yaitu :

\begin{enumerate}
\def\labelenumi{\arabic{enumi}.}
\item
  Batasan penelitian 1
\item
  Batasan penelitian 2
\item
  Batasan penelitian dapat berupa batasan ruang lingkup / sampel yang
  digunakan dalam penelitian. Contoh bentuk batasan dapat dilihat pada
  nomor-nomor selanjutnya.
\item
  Penelitian ini hanya dilakukan pada {[}nama objek penelitian{]}, tidak
  mencakup sistem atau objek lain di luar itu (Contoh saja, dapat
  disesuaikan tergantung dengan case penelitiannya).
\item
  Data penelitian diperoleh melalui {[}misalnya kuesioner dan studi
  literatur{]}, tidak mencakup metode pengumpulan data lain seperti
  wawancara mendalam (jika penelitian ini menggunakan pengumpulan data
  yang jelas, jika tidak maka tidak perlu).
\item
  Subjek penelitian dibatasi pada {[}kelompok responden tertentu,
  misalnya mahasiswa aktif angkatan 2022--2023{]} (Catatan : jika
  penelitian ini terkait untuk pengambilan data, jika tidak maka tidak
  perlu).
\item
  Variabel penelitian yang dianalisis hanya meliputi {[}sebutkan
  variabel, misalnya \emph{usability, efficiency, satisfaction}{]},
  tanpa membahas faktor lain di luar itu (Catatan : jika tipe penelitian
  lebih ke arah kuantitif menggunakan statistik, jika tidak maka tidak
  perlu).
\item
  Metode yang digunakan dalam penelitian ini adalah {[}metode yang
  dipilih, misalnya AHP dan TOPSIS{]}, tanpa membandingkan dengan metode
  lain.
\item
  Lokasi penelitian dibatasi pada {[}sebutkan lokasi, misalnya Perusahan
  XYZ{]}. Jika penelitian ini ada keinginan untuk konfidensial tanpa
  menyebutkan nama perusahaan, maka cukup bisa menyebutkan domain bidang
  perusahaan dan perusahaan dengan nama XYZ.
\item
  Data penelitian diambil dari periode {[}bulan--tahun, misalnya
  Maret--Juli 2025{]}. (Ini contoh juga bagi type penelitian yang
  menggunakan data tertentu agar tahu bahwa data yang digunakan dari
  tahun berapa sampai tahun berapa)
\item
  Hasil penelitian ini hanya berlaku sesuai dengan lingkup data dan
  responden yang digunakan, sehingga generalisasi di luar konteks
  penelitian ini memiliki keterbatasan.
\end{enumerate}

\hypertarget{manfaat-penelitian}{%
\subsection{Manfaat Penelitian}\label{manfaat-penelitian}}

Manfaat penelitian dibagi menjadi (x -- jumlah), yaitu:

\begin{enumerate}
\def\labelenumi{\arabic{enumi}.}
\item
  \begin{quote}
  Bagi Universitas Telkom, penelitian ini memberikan kontribusi untuk
  apa saja baik secara akademis maupun institusional.
  \end{quote}
\item
  \begin{quote}
  Bagi organisasi / stakeholder (sebutkan secara spesifik siapa misal
  pengelola kantor desa, perusahaan garment), dengan menjelaskan secara
  spesifik manfaat dan value (nilai dampak) yang diperoleh
  organisasi/stakeholder objek penelitian. Jangan hanya menyebutkan
  untuk siapa, tetapi tekankan "untuk apa" manfaat tersebut secara
  konkret (misal: efisiensi waktu/biaya, perbaikan proses bisnis,
  akurasi pengambilan keputusan, optimalisasi pengelolaan data,
  peningkatan kinerja SDM, memberikan \emph{insight} terkait pengambilan
  keputusan apa, dst).
  \end{quote}
\item
  \begin{quote}
  Bagi peneliti lain yang bergerak dalam bidang keilmuan yang sama,
  jelaskan penelitian ini bermanfaat untuk apa saja.
  \end{quote}
\item
  \begin{quote}
  Bagi masyarakat (sebutkan segmen masyarakat, jika ada), hasil
  penelitian ini dapat digunakan untuk apa saja.
  \end{quote}
\item
  \begin{quote}
  Bagi peneliti selanjutnya, penelitian ini digunakan untuk jelaskan apa
  kegunaan secara akademis.
  \end{quote}
\end{enumerate}

\hypertarget{sistematika-penulisan}{%
\subsection{Sistematika Penulisan}\label{sistematika-penulisan}}

Penelitian ini diuraikan dengan sistematika penulisan sebagai berikut:

\begin{longtable}[]{@{}
  >{\raggedright\arraybackslash}p{(\columnwidth - 2\tabcolsep) * \real{0.1480}}
  >{\raggedright\arraybackslash}p{(\columnwidth - 2\tabcolsep) * \real{0.8520}}@{}}
\toprule\noalign{}
\begin{minipage}[b]{\linewidth}\raggedright
\textbf{Bab I}
\end{minipage} & \begin{minipage}[b]{\linewidth}\raggedright
\textbf{Pendahuluan}

Bab ini berisi latar belakang penelitian, perumusan masalah, tujuan
penelitian, dst silakan dijelaskan
\end{minipage} \\
\midrule\noalign{}
\endhead
\bottomrule\noalign{}
\endlastfoot
\textbf{Bab II} & \textbf{KAJIAN PUSTAKA}

Bab ini membahas berbagai kajian pustaka yang relevan dengan penelitian
ini \\
\textbf{Bab III} & \textbf{Metode Penyelesaian Masalah}

Bab ini menguraikan metode yang digunakan dalam penelitian, yang terdiri
dari langkah-langkah yang akan dilakukan untuk mengatasi masalah yang
telah dirumuskan sebelumnya dst silakan diperjelas \\
& \\
& \\
& \\
\end{longtable}

\hypertarget{section}{%
\section{}\label{section}}

\hypertarget{tinjauan-pustaka}{%
\section{TINJAUAN PUSTAKA}\label{tinjauan-pustaka}}

Bab tinjauan pustaka berisi uraian tentang hasil studi pustaka yang
terkait dengan permasalahan, kerangka kerja, dan metode di topik
penelitian. Kerangka kerja, metode, dan mekanisme yang digunakan untuk
menjawab rumusan masalah harus disertakan lebih dari satu sehingga
terdapat analisa pemilihan kerangka kerja, metode, dan mekanisme yang
cocok untuk menjawab rumusan masalah. Pada bagian ini \textbf{wajib}
menggunakan \emph{reference manager} seperti \emph{\textbf{Mendeley}}
atau sejenisnya untuk membantu pembimbing dan penguji mengecek
konsistensi penggunaan referensi

\hypertarget{materi-materi-pendukung-riset-judul-sub-bab-silakan-disesuaikan-tidak-harus-seperti-ini}{%
\subsection{Materi-materi pendukung riset (Judul sub bab silakan
disesuaikan tidak harus seperti
ini)}\label{materi-materi-pendukung-riset-judul-sub-bab-silakan-disesuaikan-tidak-harus-seperti-ini}}

Ceritakan tinjauan pustaka yang mendukung penelitian ini. Pada bagian
ini tidak perlu menceritakan hal-hal yang umum, seperti misalnya HTML,
aplikasi web, aplikasi \emph{mobile}, bahasa pemrograman, dll. Jelaskan
sesuatu hal yang secara spesifik di mana belum banyak pembaca tahu yang
dapat mendukung penjelasan bagian-bagian teknis yang pembaca kurang
tahu. Misal jika ada istilah RNN (\emph{Recurent Neural Network}), pada
saat template ini dibuat, RNN termasuk hal yang baru. RNN ini dapat
diceritakan untuk mendukung pengetahuan dan penjelasan agar membantu
pembaca mengerti lebih dalam apa yang sedang ditulis.

Tinjauan pustaka dapat diturunkan dari judul penelitian atau dari
rumusan masalah. Misal jika memilih topik terkait aplikasi pengembangan
stok obat, maka dapat dimulai dengan menjelaskan terkait dengan apa itu
manajemen stok. Selanjutnya diharapkan dapat mengakitkan juga landasan
pustaka dengan kebidangan riset atau keilmuan sesuai bidang keilmuan
mahasiswa. Misal jika membahas manajemen stok, jangan terlalu banyak
dibahas dari sisi manajemennya saja tetapi juga dari sisi keilmuan
mahasiswa misal terkait dengan aplikasi, proses bisnis yang umum,
\emph{information system governance}, sisi \emph{network} dll.

Pada tugas akhir dapat mencantumkan gambar. Penulisan \emph{caption}
atau judul gambar diletakkan di bagian bawah gambar. Gambar dan
\emph{caption} dibuat rata di tengah. Pada tugas akhir ini juga dapat
mencantumkan tabel. Judul tabel diletakkan di bagian atas Tabel. Tabel
dan caption tabel adalah dibuat rata tengah juga. Usahakan tabel selalu
dalam 1 halaman. Jika lebih dari 1 halaman, pertimbangkan untuk
dipindahkan menjadi lampiran. Usahakan penomoran gambar dan tabel adalah
menggunakan \emph{tool} penomoran otomatis dari Ms. Word untuk
mempermudah pengecekan.\\
Setiap gambar dan tabel wajib disebut nomor gambarnya pada deskripsi.
Aturan penyebutan, harus disebutkan kalimatnya terlebih dahulu, baru
ditampilkan gambarnya. Bukan gambarnya dulu, baru disebutkan kalimatnya.
Contoh penyebutan gambar adalah sebagai berikut. Contoh tampilan dari
sebuah aplikasi \emph{dashboard} monitoring untuk perusahaan XYZ dapat
dilihat pada gambar II-1.

\includegraphics[width=5.60854in,height=2.72265in]{docxmedia/media/image4.png}

\protect\hypertarget{_Toc211422569}{}{}Gambar II‑1 Tampilan halaman
\emph{dashboard} monitoring

Setiap gambar harus diberi judul dengan format penomoran
nomor-romawi-Babnomor urutan-gambar dalam Bab tersebut. Judul gambar
berada di bawah gambar dan tepat di tengah kertas. Jika gambar diambil
dari suatu referensi, maka harus dicantumkan sumber gambar tersebut di
samping caption gambar. Judul atau nama gambar ditulis dengan huruf
kecil, kecuali huruf pertama kata pertama yang ditulis dengan huruf
kapital. Baris-baris judul gambar dipisahkan dengan satu spasi.

Gambar harus dapat di-update otomatis di Daftar Gambar. Untuk itu,
gunakan caption di tiap gambar dengan cara: klik kanan gambar
\textgreater{} Insert Caption. Tambahkan jenis caption ``Gambar'' jika
belum ada; jangan gunakan ``Figure'' (default bawaan di Ms. Word) jika
Tugas Akhir menggunakan Bahasa Indonesia. Jika ada perubahan isi,
halaman, atau urutan gambar, klik kanan Daftar Gambar \textgreater{}
Update Field.

Deskripsi dari gambar atau tabel harus spesifik menjelaskan informasi
yang disajikan, termasuk apabila gambar atau tabel terdapat pada
Lampiran. Setiap gambar atau tabel kemudian harus dijelaskan lebih
lanjut di paragraf isi. Gambar yang memiliki teks atau angka harus dapat
dibaca dengan baik. Jika diperlukan, gunakan layout dengan orientasi
landscape untuk gambar agar isinya dapat terbaca dengan baik.

\includegraphics[width=5.0171in,height=3.68365in]{docxmedia/media/image5.png}

\protect\hypertarget{_Toc211422570}{}{}Gambar II‑2 Perbandingan tingkat
penjualan dan jumlah produk

Sumber: PT X (2018)

Pada Gambar II-2, perbandingan antara penjualan dan jumlah produk yang
telah diproduksi. Grafik pada gambar tersebut menunjukkan volatilitas.

Contoh penyebutan tabel adalah sebagai berikut. Perbandingan sebuah
hasil pengerjaan tugas dilihat dari aspek sukses dapat dilihat seperti
pada Tabel II-1.

\protect\hypertarget{_Toc211422712}{}{}Tabel II‑1 Perbandingan contoh
tugas

\begin{longtable}[]{@{}
  >{\raggedright\arraybackslash}p{(\columnwidth - 8\tabcolsep) * \real{0.0611}}
  >{\raggedright\arraybackslash}p{(\columnwidth - 8\tabcolsep) * \real{0.4208}}
  >{\raggedright\arraybackslash}p{(\columnwidth - 8\tabcolsep) * \real{0.1786}}
  >{\raggedright\arraybackslash}p{(\columnwidth - 8\tabcolsep) * \real{0.1999}}
  >{\raggedright\arraybackslash}p{(\columnwidth - 8\tabcolsep) * \real{0.1395}}@{}}
\toprule\noalign{}
\begin{minipage}[b]{\linewidth}\raggedright
\textbf{No}
\end{minipage} & \begin{minipage}[b]{\linewidth}\raggedright
\emph{\textbf{Tugas}}
\end{minipage} & \begin{minipage}[b]{\linewidth}\raggedright
\emph{\textbf{Success Rate}}
\end{minipage} & \begin{minipage}[b]{\linewidth}\raggedright
\emph{\textbf{Misclick Rate}}
\end{minipage} & \begin{minipage}[b]{\linewidth}\raggedright
\emph{\textbf{Duration}}
\end{minipage} \\
\midrule\noalign{}
\endhead
\bottomrule\noalign{}
\endlastfoot
1 & Masuk sebagai \emph{user} dan akses halaman \emph{dashboard.} &
100\% & 22.8\% & 14.5s \\
2 & Unggah gambar tanaman cabai dan lihat hasil klasifikasi penyakit. &
94.7\% & 55\% & 51.8s \\
3 & \emph{Input} data tanaman untuk mendapatkan prediksi pertumbuhan. &
100\% & 15.4\% & 27.5s \\
4 & Konsultasi perawatan tanaman melalui chatbot. & 100\% & 26\% &
11.7s \\
\end{longtable}

Setiap tabel dalam dokumen laporan Tugas Akhir harus diacu dalam
paragraf isi. Setiap tabel harus diberi judul dengan format penomoran
nomor-urut-romawi-Babnomor-urutan tabel dalam Bab tersebut. Judul tabel
ditulis di atas tabel dengan posisi di tengah kertas. Tabel harus dapat
di-update otomatis di Daftar Tabel. Untuk itu, gunakan caption di tiap
tabel dengan cara: klik kanan tabel \textgreater{} Insert Caption.
Tambahkan jenis caption ``Tabel'' jika belum ada; jangan gunakan
``Table'' (default bawaan di Ms. Word) jika Tugas Akhir menggunakan
Bahasa Indonesia. Jika ada perubahan isi, halaman, atau urutan tabel,
klik kanan Daftar Tabel \textgreater{} Update Field.

Penulisan isi tabel menggunakan huruf Times New Roman ukuran 11 dengan
jarak spasi 1pt. Jika tabel memiliki kolom yang banyak, penulis dapat
menggunakan layout dengan orientasi landscape dengan mengakses Layout
\textgreater{} Page Setup \textgreater{} Orientation. Agar layout
berorientasi landscape tersebut tidak mempengaruhi halaman sebelum dan
sesudahnya, tambahkan section baru yang continuous dengan section
sebelumnya melalui Layout \textgreater{} Page Setup \textgreater{} Break
\textgreater{} Continuous. Jika ukuran tabel melebihi 1 halaman, maka
berikan ``table header'' atau kepala tabel di setiap halaman yang berisi
tabel. Caranya melalui klik tabel \textgreater{} Layout \textgreater{}
Repeat Header Rows.

Keterangan gambar dan tabel menggunakan caption agar penomorannya dapat
urut otomatis, sehingga daftar gambar dan tabel dapat dibangkitkan
secara otomatis. Untuk caption ini, klik kanan tabel/gambar lalu pilih
''Insert Caption''. Posisi keterangan gambar/tabel adalah ''centered''
di halaman.

Sitasi atau penggunaan referensi wajib menggunakan \emph{reference
manager} seperti \emph{Mendeley} dan sejenisnya. Format yang digunakan
adalah menggunakan APA Style. Contoh penggunaan sitasi dapat dilihat
seperti pada bagian ini (Suakanto dkk., 2021). Setiap membuat kalimat
yang bukan merupakan pernyataan pribadi, selalu sebutkan referensinya
(Hong dkk., 2023). Selalu cek dengan bagian di daftar pustaka untuk
memastikan bahwa referensi yang disitasi menggunakan \emph{Mendeley}
sudah muncul. Jangan lupa untuk generate Bibliografi pada bagian Daftar
Pustaka.

Aturan penulisan daftar pustaka menggunakan style APA dapat dilihat pada
laman
\url{https://apastyle.apa.org/instructional-aids/reference-guide.pdf}.
Kutipan tidak langsung adalah ide / konsep orang lain yang dikutip
dengan menggunakan kata-kata penulis / peneliti sendiri. Kutipan
langsung adalah ide / konsep orang lain yang disalin sesuai dengan
aslinya. Pada format APA, kutipan tidak langsung dituliskan dalam
kalimat / teks dengan mencantumkan nama pengarang dan tahun penerbitan,
tanpa menuliskan halaman karya yang dikutip, sebagai contoh: Contoh
kutipan 1: Soesanto (2018) menyatakan bahwa kolaborasi antara supplier
dan manufaktur dapat meningkatkan keefektifan perusahaan. Contoh kutipan
2: Kolaborasi antara supplier dan manufaktur dapat meningkatkan
keefektifan perusahaan (Soesanto, 2018).

Persamaan atau rumus diberi nomor diawali dengan nomor Bab dengan Angka
Romawi dan Angka Arab untuk urutan penulisan rumus. Setiap varaiabel
dari persamaan dijelaskan agar pembaca tahun variabel itu menjelaskan
apa.

\(f(x) = x + y + z + 10\) ( II‑1)

Keterangan :

x = variabel apa bisa dijelaskan

y = variabel kekuatan pengetahuan (contoh saja)

z = waktu tempuh (contoh saja)

Setiap persamaan yang sudah ditulis juga harus disebut. Sebagai contoh
penelitian ini menggunakan persamaan II-1 dan II-2 sebagai dasar dari
perhitungan kinerja.

\(f(x) = x + y + z + 10\) ( II‑2)

Ketentuan pemisahan kata sesuai dengan ketentuan tata bahasa yang
digunakan. Kata terakhir pada dasar halaman tidak boleh dipotong ke
halaman berikutnya. Penulisan istilah asing (di luar kosa kata dari
bahasa yang digunakan) harus mengikuti mengikuti cara penulisan pada
bahasa asing tersebut dan dituliskan dalam huruf miring / \emph{italic}.

Pada penulisan dokumen laporan Tugas Akhir dalam Bahasa Indonesia,
gunakan buku Pedoman Umum Ejaan Bahasa Indonesia Yang Disempurnakan,
Pedoman Umum Pembentukan Istilah, Kamus Besar Bahasa Indonesia, dan
kamus-kamus bidang khusus yang diterbitkan oleh Pusat Pembinaan dan
Pengembangan Bahasa, Departemen Pendidikan dan Kebudayaan, sebagai
pedoman. Penulisan dokumen laporan Tugas Akhir dalam Bahasa Inggris
menggunakan sistem ejaan US English (American English).

Pada subbab ini, mahasiswa diijinkan untuk membuat subbab-subbab yang
lain. Jadi bagian ini tidak terpaku dalam satu subbab saja (II.1 saja).
Setiap subbab berisi minimal dua paragraf, sehingga jika terlalu sedikit
pertimbangkan untuk dilebur. Mahasiswa dapat membuat banyak subbab
sesuai dengan keperluan menjelaskan di bagian ini.

Misalnya:\\
II. 1 Manajemen Keuangan

II.1.1 Definisi Manajemen Keuangan\\
II.1.2 Proses Bisnis dalam Manajemen Keuangan\\
II. 2. Sistam Informasi Manajemen Keuangan\\
II.2. Sistem Informasi Manajemen Keuangan\\
II.2.2 Contoh Aplikasi Manajemen Keuangan

II.3 SAP FICO

Mahasiswa dapat menuliskan bagian ini dari cerita yang umum terlebih
dahulu, baru spesifik. Bagian ini juga akan menjadi \emph{guideline}
dengan pengetahuan bagi mahasiswa ketika merancang atau membuat solusi
terkait topik penelitian. Sumber referensi atau acuan dapat dijelaskan
di sini.

Istilah dalam bahasa Inggris selalu dituliskan dalam huruf miring
sebagai contohnya misal \emph{artificial intelligence.}

\hypertarget{penelitian-terdahulu-judul-sub-bab-tidak-harus-seperti-ini-silakan-sesuaikan}{%
\subsection{Penelitian Terdahulu (Judul Sub bab tidak harus seperti ini,
silakan
sesuaikan)}\label{penelitian-terdahulu-judul-sub-bab-tidak-harus-seperti-ini-silakan-sesuaikan}}

Pada subbab ini disajikan penelitian terdahulu (tugas akhir, skripsi,
tesis, disertasi, dan lain-lain) yang terkait dengan topik permasalahan
(5 tahun terakhir). Tujuan dari bagian ini adalah untuk melihat
orisinalitas pengerjaan tugas akhir dari sisi objek kajian dan kerangka
kerja, metode, dan mekanisme yang digunakan untuk menjawab rumusan
masalah. Pilihlah penelitian terdahulu yang relevan atau terkait dengan
topik riset. Jika topik riset terkait dengan pengembangan manajemen stok
obat maka carilah penelitian (dari jurnal, prosiding) yang sejenis. Dari
situ silakan membandingkan kelebihan dan kekurangan dari masing-masing
penelitian terdahulu tersebut. Sehingga jika sudah tahu kelebihan dan
kekurangan penelitian lain, maka penelitian yang dibuat juga harus tahu
cara memposisikannya. Idealnya penelitian yang diusulkan dapat
memperbaiki atau melengkapi kekurangan atau kekosongan dari
penelitian-penelitian sebelumnya.

Sebagai ringkasan, penelitian pendahulu dapat dibuat dalam bentuk tabel.
Untuk format tabel (kolom) silakan menyesuiakan dengan kebutuhan sesuai
dengan topik penelitian. Selalu gunakan nama penulis dan tahun sebagai
bagian untuk melihat perkembangan riset di bidang tersebut dari tahun ke
tahun. Jika perbandingan pada penelitian banyak (kolomnya) banyak
silakan gunakan versi \emph{landscape} untuk lembar tersebut. Untuk
judul kolom silakan disesuaikan dengan kebutuhan penelitian. Judul per
kolom tidak harus seperti pada contoh di Tabel II-2, karena
masing-masing riset punya kekhasan yang ingin ditonjolkan atau
diperbandingkan.

Misal jika :

Peminatan EDM : Kolom pembanding bisa berisi metode, \emph{variable
input}, \emph{variable output}, hasil (akurasi, dll).

Peminatan EISD : Kolom pembanding bisa berisa metode, fitur aplikasi,
kelebihan aplikasi, kekurangan.

Peminatan SAG : Kolom pembanding bisa berupa perbandingan beberapa
metode dengan kekhasannya masing-masing.

\protect\hypertarget{_Toc211422713}{}{}Tabel II‑2 Posisi Penelitian
(\emph{State of The Art)} atau judul dapat disesuaikan dengan konteksnya

\begin{longtable}[]{@{}
  >{\raggedright\arraybackslash}p{(\columnwidth - 12\tabcolsep) * \real{0.0361}}
  >{\raggedright\arraybackslash}p{(\columnwidth - 12\tabcolsep) * \real{0.1419}}
  >{\raggedright\arraybackslash}p{(\columnwidth - 12\tabcolsep) * \real{0.0753}}
  >{\raggedright\arraybackslash}p{(\columnwidth - 12\tabcolsep) * \real{0.2494}}
  >{\raggedright\arraybackslash}p{(\columnwidth - 12\tabcolsep) * \real{0.0935}}
  >{\raggedright\arraybackslash}p{(\columnwidth - 12\tabcolsep) * \real{0.1676}}
  >{\raggedright\arraybackslash}p{(\columnwidth - 12\tabcolsep) * \real{0.2362}}@{}}
\toprule\noalign{}
\begin{minipage}[b]{\linewidth}\raggedright
\textbf{No}
\end{minipage} & \begin{minipage}[b]{\linewidth}\raggedright
\emph{\textbf{Author}}
\end{minipage} & \begin{minipage}[b]{\linewidth}\raggedright
\textbf{Metode}
\end{minipage} & \begin{minipage}[b]{\linewidth}\raggedright
\textbf{Fokus Penelitian}
\end{minipage} & \begin{minipage}[b]{\linewidth}\raggedright
\textbf{Kekuatan}
\end{minipage} & \begin{minipage}[b]{\linewidth}\raggedright
\textbf{Kekurangan}
\end{minipage} & \begin{minipage}[b]{\linewidth}\raggedright
\textbf{Kolom Pembanding yang lain}
\end{minipage} \\
\midrule\noalign{}
\endhead
\bottomrule\noalign{}
\endlastfoot
1 & (Saldan dkk., 2021) & \emph{XGBoost} & Pengembangan sistem yang
dapat mendeteksi sistem dengan cara xyz & & & \\
2 & (Rivera-Medina dkk., 2022) & & & & & \\
3 & (Morchid dkk., 2023) & & & & & \\
4 & (Kumari dkk., 2024) & & & & & \\
5 & (Tan dkk., 2024) & & & & & \\
6 & (Sharma dkk., 2024) & & & & & \\
7 & Penelitian Ini & & & & & \\
\end{longtable}

Usahakan posisi penelitian diurutkan berdasarkan urutan waktu atau
urutan tahun dari yang paling lama ke paling baru. Hal ini untuk
membantu mahasiswa melihat perkembangan penelitian sejenis/serupa dari
waktu ke waktu. Sehingga posisi penelitian saat ini (di tahun ini)
seperti apa. Apakah penelitian ini hanya mengulang atau sudah pernah
dilakukan oleh penelitian yang lain, atau penelitian ini memiliki
kontribusi yang berbeda dari penelitian yang lain? Dengan tabel ini,
maka akan mempermudah untuk melihat posisi penelitian ini dibandingkan
dengan penelitian terdahulu. Lakukan sedikit pembahasan terkait dengan
penelitian-penelitian tersebut untuk memperkuat posisi penelitian ini
dibandingkan yang lain.

Penelitian ini menunjukkan keunikan tersendiri melalui fokus pada apa,
bisa lebih dijelaskan secara lebih detail. Jelaskan kembali keunikan
atau kontribusi dari penelitian ini berdasarkan hasil dari posisi
penelitian pendahulu.

\hypertarget{metode-penyelesaian-masalah}{%
\section{METODE PENYELESAIAN
MASALAH}\label{metode-penyelesaian-masalah}}

Metode penyelesaian masalah adalah langkah-langkah yang dirancang untuk
menjawab pertanyaan dalam rumusan masalah yang telah ditetapkan
sebelumnya. Bab ini akan menguraikan tahapan penelitian secara
sistematis dan terstruktur dalam upaya menyelidiki permasalahan yang
dibahas.

\hypertarget{metode-penelitian}{%
\subsection{Metode Penelitian}\label{metode-penelitian}}

Pada bagian ini jelaskan terkait metode penelitian yang dipilih. Metode
penelitian yang dipilih ini akan menjadi \emph{guidance} dalam melakukan
penelitian. Langkah-langkah apa saja yang perlu dilakukan agar
memastikan penelitian dilakukan dengan baik. Dan setiap penelitian harus
dapat dievaluasi dengan baik juga. Pemilihan Metode penelitian yang
tepat akan memberikan \emph{insight} bagi \emph{reviewer} atau pembaca
bahwasannya penelitian telah dikerjakan dengan baik, meskipun
\emph{reviewer} tidak ikut terlibat meneliti. Akan tetapi dengan melihat
langkah-langkah penelitian sudah dijalankan dengan langkah yang tepat,
maka \emph{reviewer} juga akan memiliki pendangan bahwa penelitian telah
dijalankan dengan baik dan tidak asal-asalan atau lompat-lompat tidak
jelas.

Metodologi yang dituliskan harus mencerminkan langkah-langkah nyata yang
dilakukan oleh mahasiswa. Mahasiswa dilarang mencantumkan tahapan atau
rumus evaluasi hasil \emph{generative AI} secara mentah atau tanpa
diolah lagi yang pada kenyataannya tidak dieksekusi atau tidak dipahami
saat proses penelitian berlangsung.

Jika ada alternatif metode yang lain, jelaskan ada metode apa saja. Bisa
dijelaskan dalam bentuk tabel untuk dapat diperbandingkan antara
beberapa metode. Jelaskan kenapa metode yang digunakan untuk
menyelesaikan ini adalah metode tersebut. Pastikan bahwa metode yang
dipilih memang metode yang tepat untuk permasalahan tersebut. Pastikan
juga bahwa metode yang dipilih merupakan metode terbaru yang relevan
dalam masa kini.

\hypertarget{sistematika-penyelesaian-masalah}{%
\subsection{Sistematika Penyelesaian
Masalah}\label{sistematika-penyelesaian-masalah}}

Untuk memperjelas langkah-langkah penelitian, maka dapat digambarkan
dalam bentuk diagram/\emph{flow chart} terkait langkah-langkah
penelitian secara sistematis\emph{.} Langkah-langkah tersebut merupakan
turunan dari metode penelitian yang sudah umum dan/atau dengan
menambahkan modifikasi atau penambahan langkah-langkah yang dirasa
diperlukan. Contoh gambar sistematika penelitian dapat dilihat seperti
pada Gambar III-1. Sistematika penelitian disesuaikan dengan konteks
masing-masing bidang keilmuan, gambar III-1 yang diberikan hanya contoh
saja.

\includegraphics[width=1.98958in,height=1.97917in]{docxmedia/media/image6.jpg}

\protect\hypertarget{_Toc211422571}{}{}Gambar III‑1 Alur sistematika
penelitian

\hypertarget{tahap-pertama}{%
\subsubsection{Tahap Pertama}\label{tahap-pertama}}

Dari gambar yang dibuat, jelaskan tahap-tahapnya di bagian ini. Tahap
pertama melakukan apa saja jelaskan di sini

\hypertarget{tahap-kedua}{%
\subsubsection{Tahap Kedua}\label{tahap-kedua}}

Pada tahap kedua melakukan apa saja jelaskan di sini.

\hypertarget{metode-pengumpulan-data-dan-pengolahan-data-opsional}{%
\subsection{Metode Pengumpulan data dan Pengolahan Data
(Opsional)}\label{metode-pengumpulan-data-dan-pengolahan-data-opsional}}

Jika penelitian menggunakan data (seperti dari peminatan EDM, SAG, ERP),
jelaskan bagaimana cara mendapatkan data. Apakah mengumpulkan sendiri
dari kuisioner. Apakah mengambil dari repository, jelaskan repository
darimana? Jelaskan apakah mengambil data primer atau data sekunder.
Sumber data jelaskan ada di mana dan berikan \emph{link} untuk membantu
\emph{reviewer} melihat data-data yang sudah dikumpulkan sudah tepat
atau belum.

Jika melakukan pengolahan data, jelaskan juga seperti apa cara
pengolahan data yang dilakukan. Apakah akan mengolah data-data statistik
dengan cara seperti apa atau menggunakan \emph{tools} apa, jelaskan juga
pada bagian ini.

\hypertarget{metode-evaluasi}{%
\subsection{Metode Evaluasi}\label{metode-evaluasi}}

Jelaskan secara singkat metode evaluasi yang akan digunakan di sini.
Evaluasi dari penelitian dapat dilakukan dengan berbagai cara. Jelaskan
di sub bab ini untuk membantu para pembaca dan \emph{reviewer} untuk
mamastikan bahwa penelitian ini akan dievaluasi dengan cara yang benar
dan tepat, agar hasil yang diperoleh juga menghasilkan hasil penelitian
yang dapat dipercaya.

Metode evaluasi yang dirancang diharapkan dapat memiliki keterkaitan
langsung dengan rumusan masalah dan tujuan penelitian. Pengujian yang
dilakukan tidak boleh hanya berhenti pada pengujian teknis (seperti
Black Box Testing), melainkan harus dirancang sedemikian rupa sehingga
hasil evaluasinya dapat terukur dan secara obyektif dijadikan dasar
untuk menyimpulkan apakah permasalahan yang diangkat pada Bab I telah
berhasil diselesaikan. Sebagai contoh, jika rumusan masalah berkaitan
dengan efisiensi alur kerja atau tingkat adopsi pengguna, maka pengujian
wajib mencakup pengukuran efisiensi waktu (\emph{Process}) atau kepuasan
pengguna (\emph{People}).

Lakukan konsistensi penulisan dari bab 1 sampai dengan metode untuk
menjaga benang merah penelitian secara eksplisit, mulai dari perumusan
masalah di Bab I, perancangan evaluasi di Bab III, penyajian data di Bab
V, hingga penarikan kesimpulan di Bab VI. Apabila pengujian tidak mampu
memberikan bukti konkret untuk menjawab rumusan masalah, maka metode
evaluasi yang digunakan bisa jadi dianggap belum tepat.

Bagian evaluasi ini nanti akan dijelaskan lebih detail kembali di Bab V.
Sehingga relevansi penjelasan di subbab ini nantinya harus terhubung
dengan uraian lebih detail di Bab V.

\hypertarget{daftar-pustaka}{%
\section*{DAFTAR PUSTAKA}\label{daftar-pustaka}}
\addcontentsline{toc}{section}{DAFTAR PUSTAKA}

Abdelaziz, A. M., Alarabi, L., Basalamah, S., \& Hendawi, A. (2021). A
multi-objective optimization method for hospital admission problem---a
case study on COVID-19 patients. \emph{Algorithms}, \emph{14}(2).
https://doi.org/10.3390/a14020038

Hong, T., Cho, J., Yu, H., Ko, Y., \& Seo, J. (2023). Knowledge-grounded
dialogue modelling with dialogue-state tracking, domain tracking, and
entity extraction. \emph{Computer Speech and Language}, \emph{78}.
https://doi.org/10.1016/j.csl.2022.101460

Reffad, H., \& Alti, A. (2023). Semantic-Based Multi-Objective
Optimization for QoS and Energy Efficiency in IoT, Fog, and Cloud ERP
Using Dynamic Cooperative NSGA-II. \emph{Applied Sciences
(Switzerland)}, \emph{13}(8). https://doi.org/10.3390/app13085218

Suakanto, S., Nuryatno, E. T., Fauzi, R., Andreswari, R., \& Yosephine,
V. S. (2021). Conceptual Asset Management framework: A Grounded Theory
Perspective. \emph{2021 International Conference Advancement in Data
Science, E-Learning and Information Systems, ICADEIS 2021}.
https://doi.org/10.1109/ICADEIS52521.2021.9701948

~

\hypertarget{lampiran}{%
\section{LAMPIRAN}\label{lampiran}}

\protect\hypertarget{_Toc211427018}{}{}LAMPIRAN A JUDUL LAMPIRAN

Lampiran biasanya merupakan hal-hal tambahan yang jika dimasukkan ke
dalam bagian inti malah membuat terlalu panjang sehingga tidak enak
dibaca. Bagian yang panjang seperti gambar yang terlalu banyak, tabel
yang terlalu panjang dapat dipindahkan menjadi lampiran. Jangan lupa
selalu lampiran juga disebutkan dalam bagian inti. Misal : Hasil
wawancara kerja atau transkrip wawancara dengan user dapat dilihat pada
Lampiran A\textbf{.}

Beberapa hal yang sifatnya merupakan pendukung juga dapat dilampirkan
pada bagian lampiran.

\textbf{\hfill\break
}

\protect\hypertarget{_Toc211427019}{}{}LAMPIRAN B DOKUMENTASI WAWANCARA
(CONTOH SAJA)

\includegraphics[width=5.51181in,height=2.92847in]{docxmedia/media/image7.jpeg}

Lampiran B-. Contoh dokumentasi wawancara 1

LAMPIRAN C LAMPIRAN LAIN-LAIN

Nanti lampiran yang lain-lain bisa dilampirkan. Jangan lupa setiap
lampiran juga perlu disebutkan di bagian isi dokumen.

LAMPIRAN PENDUKUNG I : PENGGUNAAN AI UNTUK MENDUKUNG PENULISAN BAB I

AI yang digunakan dalam Bab I : ChatGPT, Gemini.

Silakan lampirkan screenshot atau dalam bentuk link google drive terkait
screenshot atau printscreen dari penggunaan AI.

Penggunaan AI Bab I dengan ChatGPT :
\href{file:///D:/01\%20Data/111\%20KK\%20DIGEST/49\%20Template\%20dan\%20panduan/2026/link\%20drive\%20google}{Google
Drive Link}

Penggunaan AI Bab I dengan Gemini :
\href{file:///D:/01\%20Data/111\%20KK\%20DIGEST/49\%20Template\%20dan\%20panduan/2026/link\%20drive\%20google}{Google
Drive Link}

Porsi Penggunaan AI

Mahasiswa diwajibkan menguraikan secara deskriptif untuk apa saja AI
tersebut digunakan (misalnya untuk membantu mengeksplorasi ide, menyusun
struktur kalimat, merapikan tata kalimat, merangkum referensi awal,
dsb). Selain penjelasan teks, mahasiswa wajib melampirkan bukti autentik
berupa tangkapan layar (\emph{screenshot}) yang diletakkan dalam tahutan
Google Drive yang memuat rekam jejak instruksi (\emph{prompt}) beserta
hasil keluaran (\emph{output}) dari masing-masing platform AI yang
digunakan, seperti ChatGPT, Gemini, atau alat bantu sejenis, untuk
memastikan penggunaan teknologi dilakukan secara etis dan bertanggung
jawab.

\textbf{\hfill\break
}

LAMPIRAN PENDUKUNG II : PENGGUNAAN AI UNTUK MENDUKUNG PENULISAN BAB II

AI yang digunakan dalam Bab II: ChatGPT, Claude.

Silakan lampirkan screenshot atau dalam bentuk link google drive terkait
screenshot atau printscreen dari penggunaan AI.

Penggunaan AI Bab II dengan ChatGPT :
\href{file:///D:/01\%20Data/111\%20KK\%20DIGEST/49\%20Template\%20dan\%20panduan/2026/link\%20drive\%20google}{Google
Drive Link}

Penggunaan AI Bab II dengan Claude:
\href{file:///D:/01\%20Data/111\%20KK\%20DIGEST/49\%20Template\%20dan\%20panduan/2026/link\%20drive\%20google}{Google
Drive Link}

Porsi Penggunaan AI

Mahasiswa diwajibkan menguraikan secara deskriptif untuk apa saja AI
tersebut digunakan (misalnya untuk membantu mengeksplorasi ide, menyusun
struktur kalimat, merapikan tata kalimat, merangkum referensi awal,
dsb). Selain penjelasan teks, mahasiswa wajib melampirkan bukti autentik
berupa tangkapan layar (\emph{screenshot}) yang diletakkan dalam tahutan
Google Drive yang memuat rekam jejak instruksi (\emph{prompt}) beserta
hasil keluaran (\emph{output}) dari masing-masing platform AI yang
digunakan, seperti ChatGPT, Gemini, atau alat bantu sejenis, untuk
memastikan penggunaan teknologi dilakukan secara etis dan bertanggung
jawab.

\textbf{\hfill\break
}

LAMPIRAN PENDUKUNG III : PENGGUNAAN AI UNTUK MENDUKUNG PENULISAN BAB III

AI yang digunakan dalam Bab III: ChatGPT, Claude.

Silakan lampirkan screenshot atau dalam bentuk link google drive terkait
screenshot atau printscreen dari penggunaan AI.

Penggunaan AI Bab III dengan ChatGPT :
\href{file:///D:/01\%20Data/111\%20KK\%20DIGEST/49\%20Template\%20dan\%20panduan/2026/link\%20drive\%20google}{Google
Drive Link}

Penggunaan AI Bab III dengan Claude:
\href{file:///D:/01\%20Data/111\%20KK\%20DIGEST/49\%20Template\%20dan\%20panduan/2026/link\%20drive\%20google}{Google
Drive Link}

Porsi Penggunaan AI

Mahasiswa diwajibkan menguraikan secara deskriptif untuk apa saja AI
tersebut digunakan (misalnya untuk membantu mengeksplorasi ide, menyusun
struktur kalimat, merapikan tata kalimat, merangkum referensi awal,
dsb). Selain penjelasan teks, mahasiswa wajib melampirkan bukti autentik
berupa tangkapan layar (\emph{screenshot}) yang diletakkan dalam tahutan
Google Drive yang memuat rekam jejak instruksi (\emph{prompt}) beserta
hasil keluaran (\emph{output}) dari masing-masing platform AI yang
digunakan, seperti ChatGPT, Gemini, atau alat bantu sejenis, untuk
memastikan penggunaan teknologi dilakukan secara etis dan bertanggung
jawab.
